\documentclass{article}

\usepackage{arxiv}

\usepackage[utf8]{inputenc} % allow utf-8 input
\usepackage[T1]{fontenc}    % use 8-bit T1 fonts
\usepackage{lmodern}        % https://github.com/rstudio/rticles/issues/343
\usepackage{hyperref}       % hyperlinks
\usepackage{url}            % simple URL typesetting
\usepackage{booktabs}       % professional-quality tables
\usepackage{amsfonts}       % blackboard math symbols
\usepackage{nicefrac}       % compact symbols for 1/2, etc.
\usepackage{microtype}      % microtypography
\usepackage{graphicx}

\title{Bayesian Small Area Estimation of Under-Five Malaria Prevalence
in West Africa}

\author{
    Godwill Zulu
   \\
    School of Mathematics, Statistics and Applied Mathematics,
University of Galway, Galway, Ireland \\
   \\
  \texttt{\href{mailto:g.zulu1@universityofgalway.ie}{\nolinkurl{g.zulu1@universityofgalway.ie}}} \\
   \And
    Stewart Hagen
   \\
    School of Mathematics, Statistics and Applied Mathematics,
University of Galway, Galway, Ireland \\
   \\
  \texttt{} \\
  }


% tightlist command for lists without linebreak
\providecommand{\tightlist}{%
  \setlength{\itemsep}{0pt}\setlength{\parskip}{0pt}}



\usepackage[numbers]{natbib}
\usepackage{booktabs}
\usepackage{longtable}
\usepackage{array}
\usepackage{multirow}
\usepackage{wrapfig}
\usepackage{float}
\usepackage{colortbl}
\usepackage{pdflscape}
\usepackage{tabu}
\usepackage{threeparttable}
\usepackage{threeparttablex}
\usepackage[normalem]{ulem}
\usepackage{makecell}
\usepackage{xcolor}
\usepackage{amsmath}
\usepackage{amssymb}
\usepackage{bm}
\usepackage{booktabs}
\usepackage{longtable}
\usepackage{array}
\usepackage{multirow}
\usepackage{wrapfig}
\usepackage{float}
\usepackage{colortbl}
\usepackage{pdflscape}
\usepackage{tabu}
\usepackage{threeparttable}
\usepackage{threeparttablex}
\usepackage[normalem]{ulem}
\usepackage{makecell}
\usepackage{xcolor}
\begin{document}
\maketitle


\begin{abstract}
Reliable district-level estimates of malaria prevalence in children
under five are essential for targeting intervention resources,
monitoring progress against the WHO Global Technical Strategy and
evaluating equity, but Demographic and Health Surveys (DHS) are designed
to be representative at the regional level rather than at finer
administrative scales. We develop a Bayesian hierarchical area-level
Fay-Herriot framework with a BYM2 spatial prior to produce smoothed
district-level prevalence estimates for under-five malaria across
Burkina Faso, Côte d'Ivoire and Ghana, using two DHS rounds per country
spanning the 2010 to 2022 period. The modelling sequence builds up from
a single-country cross-sectional baseline to country-level space-time
models, then to a three-country single-time model (Model 3) and a
three-country two-time model (Model 4) that share a single 724-district
adjacency graph including cross-border edges. A formal sensitivity
analysis assesses the cross-border adjacency assumption, the
time-invariance of covariate slopes, and the assumption of shared versus
country-specific BYM2 hyperparameters using leave-one-out
cross-validation. The cross-sectional baseline empirically prefers the
BYM2 prior over both pure ICAR and IID alternatives. The country-level
and cross-border models recover three distinct trajectories: substantial
declines in Burkina Faso and Ghana, and a marked increase in Côte
d'Ivoire between 2012 and 2021. None of three structural relaxations of
Model 4 produces a detectable improvement in expected log predictive
density, so the simpler shared-adjacency calendar-year specification is
parsimoniously adequate and is retained as the preferred specification.
\end{abstract}

\keywords{
    small area estimation
   \and
    Bayesian hierarchical model
   \and
    BYM2 spatial prior
   \and
    Fay-Herriot
   \and
    malaria prevalence
   \and
    West Africa
   \and
    Demographic and Health Surveys
  }

\section{Introduction}\label{introduction}

\label{sec:intro}

\subsection{Background}\label{background}

\subsubsection{The malaria burden globally and West
Africa}\label{the-malaria-burden-globally-and-west-africa}

Malaria remains one of the leading causes of childhood morbidity and
mortality in sub-Saharan Africa. According to the World Malaria Report
2024 \citep{who2024malaria} there were an estimated 263 million cases of
malaria and 579,000 deaths worldwide with the WHO Africa Region bearing
the most burden, accounting for at least 94\% of cases and 95\% of
deaths globally. Within the African Region, the burden falls heaviest on
children under the age of 5, who account for roughly more than 75\% of
malaria deaths, which is three quarters of all malaria deaths despite
representing only around 15 percent of the population
\citep{who2024malaria}. Plasmodium falciparum dominates regional
transmission and is sustained by the Anopheles gambiae and Anopheles
funestus species complexes that occupy the climatic envelope of West
Africa from the Sahel southwards through the Guinea zone
\citep{bhatt2015, weiss2020}. Both vectors are strongly anthropophilic
and endophilic they bite humans by preference and rest indoors after
feeding but they exploit different habitats: An. gambiae breeds in
small, sunlit, temporary pools produced by seasonal rains, while An.
funestus persists in vegetated, more permanent water bodies that smooth
the seasonal trough in transmission \citep{sinka2010dominant}.
Development of the parasite inside the mosquito is itself
temperature-sensitive, falling away below roughly 18 °C
\citep{gething2011modelling}, which is why highland and plateau zones
see lower vectorial capacity than the warmer plains around them.

These constraints translate directly into the spatial pattern of risk
this thesis aims to estimate. Elevation and rainfall together encode the
climatic suitability of a district for the vectors and the parasite; in
the Sahel the transmission season is compressed into the few months that
follow the monsoon, which is the rationale for targeting seasonal
malaria chemoprevention there rather than elsewhere
\citep{cairns2012estimating, who2012smc}. Night-time light radiance,
used here as a proxy for urbanicity, captures the well-documented
gradient by which urban districts have lower prevalence than rural ones:
built surfaces displace breeding habitat, sanitation is denser, and
human--vector contact is reduced \citep{hay2005urbanization}. Yet
climate and land cover do not act alone. Two districts with identical
environmental conditions can still differ in burden, because what
matters is how often an infectious bite reaches a susceptible child and
how quickly that child is treated. Housing quality is among the
strongest socio-economic correlates of childhood malaria children in
homes with finished walls and closed eaves have roughly half the odds of
infection of those in unimproved housing, an effect comparable to that
of an ITN \citep{tusting2017housing, tusting2019mapping} and is captured
here, alongside asset and sanitation deprivation, through the
multidimensional poverty index. Maternal literacy indexes the education
gradient in care-seeking and adherence to chemoprevention regimens,
while travel time to the nearest health facility approximates the
practical accessibility of diagnosis and treatment, which translates not
only into individual survival but, by clearing infectious reservoirs,
into the local force of infection \citep{weiss2020}. Population density
finally distinguishes the dense peri-urban districts from the sparsely
populated rural districts in which most of the West African malaria
burden is concentrated.

The case for sub-national estimation is sharpened by the way the burden
is distributed within countries. Districts only short distances apart
can differ in prevalence by tens of percentage points, reflecting local
variation in vector ecology, the length and intensity of the
transmission season, intervention coverage and socioeconomic conditions
\citep{bhatt2015, weiss2020}. A national average, even when reported
with a tight confidence interval, hides this heterogeneity entirely and
is inadequate as a basis for targeting interventions or detecting
reversals of progress before they become visible at higher aggregations.
The World Health Organization's Global Technical Strategy for Malaria
2016--2030 explicitly identifies sub-national tailoring as a strategic
priority \citep{who2023}, and the High Burden to High Impact initiative
has since made sub-national risk stratification a defining feature of
national malaria control planning.

The Demographic and Health Surveys (DHS) Programme provides the
internationally standardised population-based source of childhood
malaria infection data in sub-Saharan Africa. Since approximately 2010
the standard biomarker module has included an HRP2-based rapid
diagnostic test (RDT) for Plasmodium falciparum, administered to
children aged 6--59 months in a representative subset of sampled
households \citep{florey2014, croft2018}. The sampling design is
stratified two-stage cluster, designed to produce reliable estimates at
the regional (admin-1) level rather than at the district (admin-2)
scale. At the district level, two problems therefore arise
simultaneously. The first is incomplete coverage and the second is
sampling instability: where a district is sampled, the estimate is
typically based on a handful of clusters, so that the apparent geography
of risk in the direct estimates can be dominated by survey noise rather
than by signal. Small area estimation (SAE) is the established
statistical response to both problems. The framework borrows information
across small areas through structured priors and covariates, so that an
estimate at one area is informed not only by its own sample but also by
data from areas that resemble it \citep{rao2015}.

While a parallel tradition in malaria mapping uses unit-level
geostatistics rather than area-level SAE, such as the continuous
prevalence surfaces produced by the Malaria Atlas Project (MAP)
\citep{gething2011, bhatt2015, weiss2020}, the two approaches are
complementary. Unit-level surfaces are continuous in space and well
suited to environmental prediction, while area-level SAE produces
estimates aligned with the administrative scale at which malaria control
interventions are actually planned and delivered. For the policy
purposes that motivate this work, the area-level approach is the more
directly useful, and the MAP surfaces serve as an external benchmark to
which we return in the Discussion.

A defining methodological feature of the present work is its
multi-country structure. The substantive argument for joining three
countries on a single adjacency graph is strong: malaria transmission
ecology does not respect national borders. The study countries share the
same dominant P. falciparum parasite and the same Anopheles vectors
\citep{who2023}, the West African monsoon drives a regional
July-to-October transmission peak \citep{bhatt2015}, and seasonal labour
migration disperses the parasite across borders \citep{gething2011}. Two
adjacent districts on opposite sides of the Ghana--Burkina Faso border,
for example, experience similar rainfall and substantial bidirectional
movement, so the residual variation in prevalence after conditioning on
the dominant ecological drivers ought to be at least as smooth across
that frontier as it is across an internal district boundary
\citep{weiss2020}.

At the same time, the temporal misalignment of the survey rounds
complicates the cross-border case: the first DHS rounds in the three
countries span four calendar years (Burkina Faso 2010, Côte d'Ivoire
2012, Ghana 2014), and a spatial smoother that treats cross-border edges
in the first round as if they connected contemporaneous surfaces would
conflate spatial with temporal variation \citep{saldarriaga2020}. This
is the principal methodological tension that the modelling sequence in
this report is designed to address. Through a series of multi-country
and sensitivity specifications, we ask empirically whether accounting
for cross-border dynamics produces better inferences than within-country
alternatives, and whether the substantive conclusions depend on these
cross-border edges.

\subsection{Aims and Objectives}\label{aims-and-objectives}

\section{Methods}\label{methods}

\label{sec:methods}

\subsection{Data}\label{data}

\subsubsection{DHS data and outcome}\label{dhs-data-and-outcome}

\label{sec:dhs} The primary data source is the DHS Programme
\citep{croft2018}. We use the household member recode (PR) file from
each survey, which records every de facto household member alongside the
malaria submodule results. The outcome is the result of a rapid
diagnostic test (RDT) for Plasmodium falciparum infection administered
to children aged 6 to 59 months \citep{florey2014}. The DHS data are
confidential and access requires a project application to the DHS
Program, which the present project obtained; cluster GPS coordinates are
randomly displaced (up to 2 km in urban areas, 5 km in rural areas, with
1 per cent of rural clusters displaced up to 10 km), and no
individual-level data are reproduced in this document. Six DHS rounds
are used in this study, two per country (Table \ref{tab:surveys}). The
two rounds per country are separated by 8 to 11 calendar years;
within-round fieldwork windows span the late wet-season and early
dry-season months in each country.

\begingroup\fontsize{8}{10}\selectfont

\begin{longtable}[t]{lllr}
\caption{\label{tab:surveys}DHS surveys used in this study, by country and round. \label{tab:surveys}}\\
\toprule
Country & Survey round & Fieldwork window & Districts sampled\\
\midrule
Burkina Faso & DHS 2010 & May 2010 to Jan 2011 & 45\\
Burkina Faso & DHS 2021 & Nov 2021 to Mar 2022 & 70\\
Côte d'Ivoire & DHS 2011-12 & Dec 2011 to May 2012 & 33\\
Côte d'Ivoire & DHS 2021 & Sep 2021 to Dec 2021 & 33\\
Ghana & DHS 2014 & Sep 2014 to Dec 2014 & 207\\
\addlinespace
Ghana & DHS 2022 & Oct 2022 to Jan 2023 & 256\\
\bottomrule
\end{longtable}
\endgroup{}

\subsubsection{District panel}\label{district-panel}

\label{sec:panel} District boundaries were obtained from GADM v4.1 and
harmonised to the finest administrative geography at which malaria
control programmes are implemented in each country: admin-2 for Ghana
(260 districts), admin-3 for Burkina Faso (351 communes) and admin-3 for
Côte d'Ivoire (113 departments), giving a combined panel of N=724N = 724
N=724 districts across the three countries. Boundary changes between
rounds were resolved through majority-area cross-walks. Spatial
assignment of DHS clusters to districts uses the nearest-feature rule
rather than strict point-in-polygon, to avoid mis-assignments arising
from GPS displacement near borders.

\subsubsection{Covariates}\label{covariates}

\label{sec:covariates} Seven district-level covariates are included in
the cross-border models, selected to capture the ecological and
structural determinants of malaria transmission set out in Section
\ref{sec:intro}. Table \ref{tab:covariates} summarises each covariate
with its substantive mechanism, data source and primary reference.
Continuous raster surfaces were aggregated to the district by
population-weighted mean using WorldPop counts as the weighting surface,
and values for each district-round were extracted over the actual DHS
fieldwork months of that survey (Table \ref{tab:surveys}), so that
seasonal variation in rainfall and temperature is encoded at the data
level rather than through an explicit seasonality term. Each covariate
was standardised to zero mean and unit variance prior to model fitting.

\begingroup\fontsize{7}{9}\selectfont

\begin{longtable}[t]{l>{\raggedright\arraybackslash}p{5.5cm}ll}
\caption{\label{tab:covariates}Covariates used in the cross-border models, the substantive mechanism for the malaria-prevalence relationship, the data source and the primary reference. \label{tab:covariates}}\\
\toprule
Covariate & Mechanism & Source & Reference\\
\midrule
Elevation & Altitude lowers ambient temperature and slows parasite development inside Anopheles vectors. & SRTM 1km DEM & Bhatt et al. (2015)\\
Rainfall & Rainfall sustains vector breeding sites; wet-season totals are a standard environmental driver. & CHIRPS daily precip. & Funk et al. (2015)\\
Temperature & Temperature governs vector survival and the extrinsic incubation period; optimum near 25 C. & MODIS LST & Bhatt et al. (2015)\\
Nighttime lights & Urbanisation reduces larval habitat and improves housing and healthcare access. & VIIRS composite & Elvidge et al. (2021)\\
Travel time & Accessibility from major centres reflects the practical reach of interventions and treatment. & Friction surface & Weiss et al. (2020)\\
\addlinespace
Pop. density & Density modifies effective transmission intensity and human-vector contact rates. & WorldPop + shapefile & WHO (2023)\\
Poverty & Household-level deprivation drives ITN ownership, care-seeking and housing quality. & Census, district (GHA), region (BFA, CIV) & Florey (2014)\\
\bottomrule
\end{longtable}
\endgroup{}

\paragraph{Elevation (SRTM).}

Altitude lowers ambient temperature and slows the sporogonic development
of Plasmodium falciparum inside the mosquito
\citep{gething2011modelling, bhatt2015}. Taken from the SRTM 1 km
digital elevation model and treated as time-invariant over the study
window.

\paragraph{Rainfall (CHIRPS).}

Rainfall sustains the temporary, sunlit pools in which Anopheles gambiae
s.l. breeds and shapes the regional north-to-south transmission gradient
\citep{sinka2010dominant, bhatt2015}. Taken from the CHIRPS daily
product at 0.05° resolution \citep{funk2015chirps} and averaged over the
DHS fieldwork months of each round.

\paragraph{Land surface temperature (MODIS LST).}

Temperature governs adult vector survival and the duration of the
extrinsic incubation period, with effective transmission falling away
below roughly 18 °C \citep{gething2011modelling}. Taken from the MODIS
LST 8-day composite at 1 km resolution and averaged over the DHS
fieldwork months.

\paragraph{Nighttime light radiance (VIIRS).}

Nighttime lights proxy urbanicity, capturing the well-documented
urban--rural prevalence gradient driven by reduced larval habitat,
denser infrastructure and lower indoor human--vector contact
\citep{hay2005urbanization, tusting2019mapping}. Taken from the VIIRS
Day/Night Band annual composite \citep{elvidge2021} for each survey year
and log-transformed prior to standardisation.

\paragraph{Travel time to the nearest health facility.}

Travel time approximates the practical accessibility of diagnosis and
treatment, which translates not only into individual survival but, by
clearing infectious reservoirs, into the local force of infection
\citep{weiss2020}. Taken from the global friction-surface travel-time
map of \citep{weiss2020}.

\paragraph{Population density (WorldPop).}

Population density modifies human--vector contact rates and complements
the nightlights covariate by separating the crowding component from the
urbanisation component. Computed from the WorldPop gridded 100-m
population estimates \citep{tatem2017worldpop} divided by district land
area and log-transformed prior to standardisation.

\paragraph{Multidimensional poverty index.}

MPI indexes the deprivation gradient that drives ITN ownership,
care-seeking, housing quality and IRS acceptance
\citep{tusting2017housing, tusting2019mapping}. For Ghana, the
district-level MPI was extracted from the Ghana Statistical Service
StatsBank API \citep{florey2014}; for Burkina Faso (RGPH 2019) and Côte
d'Ivoire (ENV 2015), MPI is published only at admin-1 level and the
regional value was assigned uniformly to every district within the
region, so the MPI slope for these two countries is identified by
between-region variation. This regional-imputation caveat is revisited
in Section \ref{sec:limitations}.

\subsection{Direct estimation}\label{direct-estimation}

\label{sec:direct} For each sampled district ii i in survey round tt t,
the design-based direct estimate of malaria prevalence is the Hájek
(weighted ratio) estimator

\begin{equation}

\hat{p}{it} ;=; \frac{\sum{k \in S_{it}} w_k , m_k}{\sum_{k \in S_{it}} w_k},

\label{eq:hajek}

\end{equation}

where SitS\_\{it\} Sit\hspace{0pt} is the set of children sampled in
district ii i at round tt t, wkw\_k wk\hspace{0pt} is the survey weight
of child kk k accounting for cluster and household selection
probabilities, and mk∈\{0,1\}m\_k \in \{0, 1\} mk\hspace{0pt}∈\{0,1\} is
the child's RDT result. The Hájek form is preferred over the
Horvitz--Thompson estimator because the denominator is itself estimated
from the survey, which reduces variability when within-district sample
sizes are small \citep{cochran1977, lumley2010}. The design-based
sampling variance
Var\textsuperscript{(p}it)\widehat{\mathrm{Var}}(\hat{p}\_\{it\})
Var(p\^{}\hspace{0pt}it\hspace{0pt}) is computed by Taylor-series
linearisation around the ratio \citep{binder1983}, with stratum and
primary-sampling-unit information taken from the DHS recode files and
supplied to a \texttt{svydesign} specification in R \citep{lumley2004}.
Because malaria prevalence is bounded on the unit interval and the
Fay--Herriot likelihood in Section \ref{sec:fh} requires a quantity on
an unbounded real line, both the direct estimate and its variance are
transferred to the logit scale. The logit-transformed estimate is

\begin{equation}

y_{it} ;=; \mathrm{logit}(\hat{p}{it}),

\label{eq:logit}

\end{equation}

and the logit-scale sampling variance follows from the delta method,

\begin{equation}

\psi{it}^{2} ;=; \frac{\widehat{\mathrm{Var}}(\hat{p}{it})}{\bigl[\hat{p}{it}(1 - \hat{p}_{it})\bigr]^{2}},

\label{eq:psi}

\end{equation}

following standard practice in area-level malaria and mortality SAE
\citep{mercer2015, wakefield2020}. Direct estimates of exactly zero or
one are handled by a small continuity adjustment before the logit
transform; districts with zero observed clusters in a round are flagged
as unobserved and enter the model only through the spatial and covariate
structure. The pair (yit,ψit2)(y\_\{it\}, \psi\emph{\{it\}\^{}\{2\})
(yit\hspace{0pt},ψit2\hspace{0pt}) summarises everything the DHS round
contributes about district ii i on the logit scale: yity}\{it\}
yit\hspace{0pt} is the design-based point estimate,
ψit2\psi\_\{it\}\^{}\{2\} ψit2\hspace{0pt} its design-based variance
treated as known by the modelling step that follows. The Hájek estimator
and its design-based variance are well behaved when the per-district
sample size is large and the design effect is moderate; at the district
scale of these DHS surveys neither condition holds reliably, which
motivates the model-based approach that follows. The empirical
sample-size, coefficient-of-variation and aggregation patterns
underlying this concern are presented in Section \ref{sec:results}.
Statistical framework \label{sec:framework} The Fay--Herriot framework
\label{sec:fh} We work within the area-level small area estimation
framework of \citep{fay1979}. The framework has two components: a
sampling model that links the design-based direct estimate to an
unobserved latent quantity, and a linking model that specifies how the
latent quantity depends on covariates and on a spatially structured
residual. The sampling model is

\begin{equation}

y_{it} \mid \theta_{it} ;\sim; \mathcal{N}!\left( \theta_{it}, ; \psi_{it}^{2} \right),

\label{eq:fh-lik}

\end{equation}

with components

\begin{itemize}

\item yit=logit(p^it)y_{it} = \mathrm{logit}(\hat{p}_{it})
yit​=logit(p^​it​), the logit-scale direct estimate from Section \ref{sec:direct};

\item θit=logit(pit)\theta_{it} = \mathrm{logit}(p_{it})
θit​=logit(pit​), the latent logit-scale true prevalence (the unknown of interest);

\item ψit2\psi_{it}^{2}
ψit2​, the design-based logit-scale sampling variance from equation \ref{eq:psi}, treated as known.

\end{itemize}

Treating ψit2\psi\emph{\{it\}\^{}\{2\} ψit2\hspace{0pt} as known is the
defining modelling choice of the Fay--Herriot approach: the sampling
variance is a property of the survey design (cluster count, weights,
within-cluster homogeneity) and is encoded in the design metadata
supplied with the recode files. With both ψit2\psi}\{it\}\^{}\{2\}
ψit2\hspace{0pt} and the model-based residual variance free, the
parameters would be only weakly identified \citep{you2002, rao2015}. We
adopt the treated-as-known assumption throughout, including for
lightly-sampled districts; smoothing ψit2\psi\emph{\{it\}\^{}\{2\}
ψit2\hspace{0pt} via a generalised variance function \citep{rao2015} is
a possible refinement we do not pursue here. The linking model specifies
the structure of the latent surface θit\theta}\{it\} θit\hspace{0pt}. It
takes the generic form

\begin{equation}

\theta_{it} \;=\; \mu + \mathbf{x}_{it}^{\top} \bm{\beta} \;+\; b_i \;+\; (\text{time, country, interaction terms}),

\label{eq:fh-link}

\end{equation}

with components

\begin{itemize}

\item μ\mu
μ, a global intercept on the logit scale;

\item xit\mathbf{x}_{it}
xit​, the seven-vector of standardised covariates for district ii
i in round tt
t, and β\bm{\beta}
β the corresponding slope vector;

\item bib_i
bi​, a district-level spatial random effect specified in Section \ref{sec:spatial};

\item time, country and interaction terms that are zero in the cross-sectional baseline and that build up across the model sequence in Section \ref{sec:models}.

\end{itemize}

The full breakdown of the linear predictor is the point at which the
model sequence branches: the cross-sectional baseline (Model 1) sets the
time and country terms to zero and probes three alternative forms of
bib\_i bi\hspace{0pt}, while the subsequent multi-country specifications
add country, time and interaction structure on top of a fixed BYM2 form
for bib\_i bi\hspace{0pt}. Spatial structure: the adjacency graph, ICAR
and BYM2 priors \label{sec:spatial} The district-level spatial random
effect bib\_i bi\hspace{0pt} is specified through a prior on the joint
distribution of (b1,\ldots,bN)(b\_1, \ldots, b\_N)
(b1\hspace{0pt},\ldots,bN\hspace{0pt}) that encodes the substantive
expectation of smoothness between neighbouring districts. The structure
we use throughout the multi-country specifications is the BYM2
reparameterisation of \citep{riebler2016}, which sits on an adjacency
graph that we describe first.

\paragraph{The adjacency graph.}

A district-level adjacency graph encodes the qualitative neighbour
relation across the panel of N=724N = 724 N=724 districts. We construct
three variants. The within-country graph WGHAW\_\{\text{GHA}\}
WGHA\hspace{0pt} over the 260 Ghanaian districts is used by the Model 1
cross-sectional baseline. The combined within-country graph stacks the
three country-only graphs on the same 724-node set, producing three
disconnected components, and is used by the country-specific BYM2
sensitivity (Model 4b) and by the edge-deletion sensitivity. The
principal multi-country specification (Models 3 and 4) uses the
cross-border graph WW W, which extends the combined graph by adding
edges between districts on opposite sides of an international border
that share a land frontier. Cross-border edges were identified manually
from harmonised admin-2/admin-3 polygons. The cross-border graph WW W
has E=2,056E = 2\{,\}056 E=2,056 edges in total: 1,978 within-country
edges and 78 cross-border edges. The cross-border edges are the
structural feature that makes this work a multi-country specification
rather than a stack of three independent country models. Their
substantive motivation is set out in Section \ref{sec:intro}: malaria
transmission ecology does not respect political boundaries, the West
African monsoon drives a regional wet-season transmission peak that
propagates across the three countries \citep{bhatt2015}, and seasonal
labour migration disperses parasites across borders \citep{gething2011}.
Two districts that face each other across a national border share
rainfall, temperature, vector ecology and human movement, so the
residual variation in prevalence after adjustment for covariates ought
to be at least as smooth across the border as it is across an internal
district boundary \citep{weiss2020}. Whether the cross-border edges
actually buy predictive performance in practice is the empirical
question that the edge-deletion sensitivity (Section \ref{sec:models})
addresses.

\paragraph{ICAR prior.}

The intrinsic conditional autoregressive (ICAR) prior of
\citep{besag1991} is the canonical prior for a smooth spatial field on a
graph. For a vector u=(u1,\ldots,uN)⊤u = (u\_1, \ldots,
u\_N)\^{}\{\top\} u=(u1\hspace{0pt},\ldots,uN\hspace{0pt})⊤ on the node
set of an adjacency matrix WW W, the ICAR density is

\begin{equation}

p(u \mid W) \;\propto\; \exp\!\left( -\tfrac{1}{2} \sum_{(i, j) \in E(W)} (u_i - u_j)^{2} \right),

\label{eq:icar}

\end{equation}

with components

\begin{itemize}

\item WW
W, the N×NN \times N
N×N adjacency matrix with Wij=1W_{ij} = 1
Wij​=1 if districts ii
i and jj
j are neighbours and 00
0 otherwise;

\item E(W)E(W)
E(W), the set of edges of WW
W (i.e., the pairs (i,j)(i, j)
(i,j) with Wij=1W_{ij} = 1
Wij​=1);

\item uiu_i
ui​, the structured spatial random effect at district ii
i.

\end{itemize}

The ICAR density is improper because it is invariant to a constant shift
in uu u; identifiability is restored by the sum-to-zero constraint
∑iui=0\sum\_i u\_i = 0 ∑i\hspace{0pt}ui\hspace{0pt}=0. For graphs with
multiple disconnected components, such as the combined within-country
graph, the sum-to-zero constraint is applied per component
\citep{freni2018}.

\paragraph{BYM2 prior.}

The pure ICAR specification forces all residual spatial variation to be
structured; the pure IID specification forces it all to be unstructured.
The BYM2 reparameterisation of \citep{riebler2016} combines the two into
a single mixture with an interpretable marginal scale,

\begin{equation}

b_i ;=; \sigma_b \left( \sqrt{1 - \phi} ; v_i + \sqrt{\phi / s} ; u_i \right),

\label{eq:bym2}

\end{equation}

with components

\begin{itemize}

\item σb>0\sigma_b > 0
σb​>0, the marginal standard deviation of the spatial random effect on the logit scale;

\item ϕ∈[0,1]\phi \in [0, 1]
ϕ∈[0,1], the mixing parameter that governs the balance between structured and unstructured variation: ϕ=0\phi = 0
ϕ=0 recovers a pure IID prior, ϕ=1\phi = 1
ϕ=1 recovers a pure ICAR prior, and intermediate values are a mixture;

\item vi∼N(0,1)v_i \sim \mathcal{N}(0, 1)
vi​∼N(0,1) independently across districts, the unstructured component;

\item ui∼ICAR(W)u_i \sim \mathrm{ICAR}(W)
ui​∼ICAR(W) as in equation \ref{eq:icar}, the structured component;

\item ss
s, a scaling factor computed from the generalised inverse of the ICAR structure matrix so that σb\sigma_b
σb​ has a consistent marginal-SD interpretation regardless of the adjacency structure [@riebler2016]. For graphs with disconnected components the scaling factor is computed per component [@freni2018; @donegan2021].

\end{itemize}

Treating ϕ\phi ϕ as an unknown parameter lets the data determine the
mixture rather than imposing one or the other \citep{morris2019}. The
BYM2 prior is the standard form used in Models 2 through 4 below; the
Model 1 cross-sectional baseline compares it to its two limiting cases
(pure IID and pure ICAR). Model sequence \label{sec:models} The notation
is shared throughout: ii i indexes districts, t∈\{1,2\}t \in \{1, 2\}
t∈\{1,2\} indexes survey rounds, c∈\{1,2,3\}c \in \{1, 2, 3\}
c∈\{1,2,3\} indexes countries (c=1c = 1 c=1 Ghana, c=2c = 2 c=2 Côte
d'Ivoire, c=3c = 3 c=3 Burkina Faso), and kk k indexes the seven
covariates. The notation c{[}i{]}c{[}i{]} c{[}i{]} denotes the country
containing district ii i. All models share the Fay--Herriot sampling
model (equation \ref{eq:fh-lik}); they differ only in the linear
predictor for θit\theta\_\{it\} θit\hspace{0pt} and, in the
cross-sectional baseline, in the form of the spatial random effect.
Model 1: Cross-sectional Ghana baseline \label{sec:m1} Model 1 fits the
Fay--Herriot framework to Ghana 2022 alone (a single country, a single
round), in three variants that share the same likelihood and the same
covariate component but differ only in the form of the district-level
spatial random effect. The shared linear predictor is

\begin{equation}

\theta_i ;=; \mu + \mathbf{x}_i^{\top} \bm{\beta} + b_i,

\label{eq:m1}

\end{equation}

with components

\begin{itemize}

\item μ\mu
μ, the Ghana 2022 grand mean on the logit scale;

\item xi\mathbf{x}_i
xi​, the standardised seven-covariate vector for district ii
i at the 2022 round, and β\bm{\beta}
β the corresponding seven-vector of slopes;

\item bib_i
bi​, the district-level spatial residual, whose three competing forms define the variants.

\end{itemize}

The covariate component xi⊤β\mathbf{x}\_i\^{}\{\top\} \bm{\beta}
xi⊤\hspace{0pt}β is identical across the three variants; only the
residual structure bib\_i bi\hspace{0pt} changes.

\paragraph{IID variant.}

The residual is fully unstructured noise,

\begin{equation}

b_i ;\sim; \mathcal{N}(0, \sigma_b^{2}),

\label{eq:m1-iid}

\end{equation}

with σb\textgreater0\sigma\_b \textgreater{} 0
σb\hspace{0pt}\textgreater0 the residual standard deviation. Two
neighbouring districts share no information beyond the covariates.

\paragraph{Pure ICAR variant.}

The residual is fully spatially structured,

\begin{equation}

b_i ;\sim; \mathrm{ICAR}(W_{\text{GHA}}),

\label{eq:m1-icar}

\end{equation}

with WGHAW\_\{\text{GHA}\} WGHA\hspace{0pt} the within-Ghana adjacency
matrix over the 260 Ghanaian districts and the sum-to-zero
identifiability constraint as in equation \ref{eq:icar}. Information
flows between neighbours but unstructured residual variation at the
district scale is forced to zero.

\paragraph{BYM2 variant.}

The residual is the mixture form of equation \ref{eq:bym2},

\begin{equation}

b_i ;=; \sigma_b \left( \sqrt{1 - \phi} ; v_i + \sqrt{\phi / s_{\text{GHA}}} ; u_i \right),

\label{eq:m1-bym2}

\end{equation}

with ui∼ICAR(WGHA)u\_i \sim \mathrm{ICAR}(W\_\{\text{GHA}\})
ui\hspace{0pt}∼ICAR(WGHA\hspace{0pt}), vi∼N(0,1)v\_i \sim \mathcal{N}(0,
1) vi\hspace{0pt}∼N(0,1), and sGHAs\_\{\text{GHA}\} sGHA\hspace{0pt} the
scaling factor computed on WGHAW\_\{\text{GHA}\} WGHA\hspace{0pt}. The
mixing parameter ϕ\phi ϕ is estimated from the data. The substantive
purpose of Model 1 is to determine empirically which form of the spatial
residual the data support, before fixing the form for the multi-country
specifications that follow. The three variants are compared in Section
\ref{sec:results}. Model 2: Single-country space--time \label{sec:m2}
Model 2 extends Model 1 to both DHS rounds within a single country and
is fitted independently for each of the three countries. For country cc
c, the linear predictor is

\begin{equation}

\theta_{it} \;=\; \mu_c + \mathbf{x}{it}^{\top} \bm{\beta}c + b_i + \delta_c , \widetilde{\text{yr}}{it} + \gamma{it},

\label{eq:m2}

\end{equation}

with components

\begin{itemize}

\item μc\mu_c
μc​, the country-cc
c intercept on the logit scale, identified as the prevalence at yr~=0\widetilde{\text{yr}} = 0
yr​=0;

\item xit\mathbf{x}_{it}
xit​, the seven-covariate vector for district ii
i at round tt
t (now time-varying for rainfall, temperature, nighttime lights and population density), and βc\bm{\beta}_c
βc​ the country-cc
c slope vector;

\item bib_i
bi​, the BYM2 spatial residual on the within-country adjacency WcW_c
Wc​ (equation \ref{eq:bym2}), shared across both rounds for district ii
i;

\item yr~it=yrit−2016\widetilde{\text{yr}}_{it} = \text{yr}_{it} - 2016
yr​it​=yrit​−2016, the calendar year of survey tt
t in district ii
i centred at the regional midpoint 2016;

\item δc\delta_c
δc​, the country-cc
c annual rate of change in logit prevalence;

\item γit∼N(0,σγ2)\gamma_{it} \sim \mathcal{N}(0, \sigma_\gamma^{2})
γit​∼N(0,σγ2​), an unstructured space–time interaction following the Type I classification of [@knorrheld2000].

\end{itemize}

Fitting Model 2 separately for each country provides the within-country
benchmark against which the multi-country specifications are compared.
Model 3: Three-country single-time \label{sec:m3} Model 3 is the first
multi-country specification. It joins all three countries on the
cross-border adjacency graph WW W from Section \ref{sec:spatial}, but
uses only the temporally aligned second survey round (Burkina Faso 2021,
Côte d'Ivoire 2021, Ghana 2022). With every district surveyed within a
one-year window, cross-border smoothing operates on contemporaneous
surfaces by construction. The linear predictor is

\begin{equation}

\theta_i ;=; \mu + \alpha_{c[i]} + \mathbf{x}i^{\top} \bm{\beta}{c[i]} + b_i,

\label{eq:m3}

\end{equation}

with components

\begin{itemize}

\item μ\mu
μ, the regional grand mean on the logit scale;

\item αc\alpha_{c}
αc​, the country-cc
c deviation from the regional mean, identified by the sum-to-zero constraint ∑cαc=0\sum_c \alpha_c = 0
∑c​αc​=0;

\item xi\mathbf{x}_i
xi​ and βc\bm{\beta}_c
βc​, district-level covariates and country-specific slopes; the slopes are partially pooled across countries through the hierarchical prior βc,k∼N(βˉk,τk2)\beta_{c, k} \sim \mathcal{N}(\bar{\beta}_k, \tau_k^{2})
βc,k​∼N(βˉ​k​,τk2​) [@gelman2006];

\item bib_i
bi​, the BYM2 spatial residual on the cross-border graph WW
W (equation \ref{eq:bym2}), with shared hyperparameters (σb,ϕ)(\sigma_b, \phi)
(σb​,ϕ) across the three countries.

\end{itemize}

The shared BYM2 hyperparameters and the cross-border edges of WW W
together encode the multi-country structure: information flows between
districts on opposite sides of an international border, and the marginal
scale of residual spatial variation is assumed common across the three
countries. Model 4: Three-country two-time \label{sec:m4} Model 4
extends Model 3 to both DHS rounds. The temporal effect enters as a
country-specific linear slope on the centred calendar year:

\begin{equation}

\theta_{it} ;=; \mu + \alpha_{c[i]} + \mathbf{x}{it}^{\top} \bm{\beta}{c[i]} + b_i + \delta_{c[i]} , \widetilde{\text{yr}}{it} + \gamma{it},

\label{eq:m4}

\end{equation}

with components

\begin{itemize}

\item μ,αc,βc\mu, \alpha_c, \bm{\beta}_c
μ,αc​,βc​ as in Model 3;

\item bib_i
bi​, BYM2 on the cross-border graph WW
W, time-invariant within a district (a single spatial surface is assumed to describe the residual spatial gradient across both rounds);

\item yr~it=yrit−2016\widetilde{\text{yr}}_{it} = \text{yr}_{it} - 2016
yr​it​=yrit​−2016, the centred calendar year of survey tt
t in district ii
i;

\item δc\delta_c
δc​, the country-cc
c annual rate of change in logit prevalence, with the country-level hierarchical prior δc∼N(δˉ,σδ2)\delta_c \sim \mathcal{N}(\bar{\delta}, \sigma_\delta^{2})
δc​∼N(δˉ,σδ2​);

\item γit∼N(0,σγ2)\gamma_{it} \sim \mathcal{N}(0, \sigma_\gamma^{2})
γit​∼N(0,σγ2​), unstructured space–time interaction (Type I, [@knorrheld2000]).

\end{itemize}

Two notes on parameterisation. First, the temporal effect is expressed
on calendar-year scale rather than as a round-to-round wave indicator
because the surveys are not synchronised across countries (Round 1 spans
2010 to 2014, Round 2 spans 2021 to 2022): a wave indicator would
conflate the country-specific elapsed time between rounds with the
country-specific rate of change \citep{wakefield2020}. The between-round
odds ratio for country cc c is recovered as exp⁡(δc⋅Δyrc)\exp(\delta\_c
\cdot \Delta \text{yr}\_c) exp(δc\hspace{0pt}⋅Δyrc\hspace{0pt}) where
Δyrc\Delta \text{yr}\emph{c Δyrc\hspace{0pt} is the actual elapsed years
between rounds. Second, bib\_i bi\hspace{0pt} is time-invariant rather
than time-varying: the dominant structural drivers of residual spatial
variation (long-run climate, elevation, the durable component of
urbanisation, accessibility) move on decadal timescales
\citep{bhatt2015, weiss2020}, and with only two rounds per district a
per-round spatial field bi,tb}\{i, t\} bi,t\hspace{0pt} is not
separately identified from a saturated round-by-district interaction.
Sensitivity specifications \label{sec:sens} Three sensitivity variants
probe the structural assumptions of Model 4.

\paragraph{Model 4a (country-year covariate slopes).}

Model 4a relaxes the time-invariance of the covariate slopes by allowing
the covariate--prevalence relationship to differ between rounds within
each country,

\begin{equation}

\theta_{it} ;=; \mu + \alpha_{c[i]} + \mathbf{x}{it}^{\top} \bm{\beta}{c[i], t} + b_i + \delta_{c[i]} , \widetilde{\text{yr}}{it} + \gamma{it},

\label{eq:m4a}

\end{equation}

with βc,t,k∼N(βˉk,τk2)\beta\_\{c, t, k\}
\sim \mathcal{N}(\bar\{\beta\}\_k, \tau\_k\^{}\{2\})
βc,t,k\hspace{0pt}∼N(βˉ\hspace{0pt}k\hspace{0pt},τk2\hspace{0pt}). This
is the country-by-year interaction specification: within each country
the covariate slopes are free to differ by round; across the six
country-rounds the slopes are partially pooled toward a single regional
mean per covariate. With only two rounds per district, country-by-year
is the finest interaction supportable here.

\paragraph{Model 4b (country-specific BYM2).}

Model 4b replaces the shared BYM2 hyperparameters of Model 4 with
country-specific ones \citep{donegan2021}. The cross-border edges are
removed from the adjacency graph (the model reverts to the disconnected
combined graph), and each country gets its own σb,c\sigma\_\{b, c\}
σb,c\hspace{0pt} and ϕc\phi\_c ϕc\hspace{0pt}:

\begin{equation}

b_i \;=\; \sigma_{b, c[i]} \left( \sqrt{1 - \phi_{c[i]}} \; v_i + \sqrt{\phi_{c[i]} / s_{c[i]}} \; u_i \right),

\label{eq:m4b}

\end{equation}

with uu u ICAR on the within-country adjacency W′W' W′, per-country
sum-to-zero ∑i∈cui=0\sum\_\{i \in c\} u\_i = 0
∑i∈c\hspace{0pt}ui\hspace{0pt}=0, and per-country scaling factor scs\_c
sc\hspace{0pt}. The substantive question is whether the three countries
have materially different residual spatial structures once the
cross-border edges are removed.

\paragraph{Edge-deletion sensitivity.}

The edge-deletion sensitivity isolates the contribution of the 78
cross-border edges by removing them and refitting Model 4 with all other
features (covariate slopes, country intercepts, temporal slopes, shared
BYM2 hyperparameters) unchanged. The scaling factor is computed per
component and applied per district \citep{freni2018}. Comparing the fit
with and without the cross-border edges directly quantifies what
cross-border smoothing buys, both for predictive performance and for the
substantive parameter estimates (country slopes, regional covariate
effects). Priors \label{sec:priors} Weakly informative priors are used
throughout, with the priors on the spatial standard deviation following
the penalised-complexity (PC) framework of \citep{simpson2017}. The
shared prior specification across the cross-border models is given in
Table \ref{tab:priors}.

\begin{table}[!h]
\centering
\caption{\label{tab:priors}Prior distributions for the shared parameters across the cross-border models. \label{tab:priors}}
\centering
\fontsize{8}{10}\selectfont
\begin{tabular}[t]{llll}
\toprule
Parameter & Prior & Rationale & Source\\
\midrule
$\mu$ & $\mathcal{N}(-1.85,\, 1.5^2)$ & Centred on logit of moderate paediatric prevalence & Gelman (2006)\\
$\alpha_c$ (free) & $\mathcal{N}(0,\, 1.5^2)$ & Identifies free elements under sum-to-zero & Gelman (2006)\\
$\sigma_b$ & $\mathrm{Exp}(4.6)$ & PC prior, $P(\sigma_b > 1) \approx 0.01$ & Simpson et al. (2017)\\
$\phi$ & $\mathrm{Beta}(3,\, 2)$ & Mild preference for structured component & Morris et al. (2019)\\
$\sigma_\gamma$ & $\mathrm{Exp}(2)$ & PC prior on space-time interaction SD & Simpson et al. (2017)\\
\addlinespace
$\bar\beta_k$ & $\mathcal{N}(0,\, 1)$ & Weakly informative on standardised covariates & Gelman (2006)\\
$\tau_k$ & $\mathrm{Exp}(2)$ & PC prior on between-country slope SD & Simpson et al. (2017)\\
$\bar\delta$ & $\mathcal{N}(0,\, 0.5^2)$ & Tight on regional annual rate of change & Wakefield et al. (2020)\\
$\sigma_\delta$ & $\mathrm{Exp}(5)$ & PC prior on between-country slope variation & Simpson et al. (2017)\\
\bottomrule
\end{tabular}
\end{table}

\#\#Computation and diagnostics

\label{sec:compute} All models were fitted in Stan version 2.34
\citep{stan2024} via the cmdstanr interface, using the No-U-Turn Sampler
\citep{hoffman2014}. Four chains were run in parallel with 1,500 warmup
and 2,000 sampling iterations each. Cross-sectional and country-level
models used adapt\_delta = 0.99 and max\_treedepth = 14; cross-border
models used adapt\_delta = 0.97 and max\_treedepth = 13. Convergence was
assessed by the rank-normalised R\^{}\hat{R} R\^{} of
\citep{vehtari2021}, bulk and tail effective sample sizes, the number of
divergent transitions reported by the sampler, and visual inspection of
trace, rank and pairs plots. Non-centred parameterisations were used for
all hierarchical components to mitigate funnel geometry
\citep{betancourt2015}.

\bibliographystyle{unsrtnat}
\bibliography{references.bib}


\end{document}
